% GENERATED FILE. Edit canonical lessons, then run npm run generate:books.
% canonical-source-hash: fnv1a64:639fc6b334532660
% canonical-lessons: TA-C58-cow, TA-C58-goat, TA-C58-hen, TA-C58-crow, TA-C58-bird

\chapter{Around the Yard}
\label{ch:ta-around-the-yard}
\hlchaptermodality{58}

\begin{chapteropening}
\textbf{By the end of this chapter:} \emph{I can name a cow, a goat, a hen, a crow and a bird in Tamil.}

Complete TA-C58-bird and say all five words in order.
\end{chapteropening}

\section[\texorpdfstring{\ta{பசு} (pasu)}{பசு (pasu)}]{\ta{பசு} (pasu) — a cow}

\label{lesson:TA-C58-cow}



\begin{quote}
Before the new one: what did \emph{pukai} mean, and what did \emph{kaṉal} mean?
\end{quote}

\subsection*{You'll want to know: \ta{பசு}}
\textbf{\ta{பசு}} (\emph{pasu}) — \textquotedblleft{}a cow\textquotedblright{}.

The animal a Tamil household is built around. Kannada says \ka{ಹಸು} (\emph{hasu}): the same word, with the \emph{p}-for-\emph{h} trade Kannada makes right through its old words — the one this book pointed out when the numbers were lined up side by side. Tamil keeps the \emph{p}.

It belongs to a \ta{விவசாயி} before it belongs to anybody, and in a \ta{கிராமம்} it stands tied in the yard rather than shut in a shed, near enough to the \ta{வாசல்} that you step around it going in.

\textbf{In the family, and next door}

\noindent
\begin{tabularx}{\linewidth}{@{}>{\raggedright\arraybackslash}X>{\raggedright\arraybackslash}X>{\raggedright\arraybackslash}X@{}}
\toprule
\textbf{} & \textbf{word} & \textbf{same root} \\
\midrule
Kannada & \ka{ಹಸು} (\emph{hasu}) & yes \\
Telugu & \te{ఆవు} (\emph{āvu}) &  \\
Malayalam & \ml{പശു} (\emph{paśu}) & yes \\
Hindi (neighbour) & \dv{गाय} (\emph{gāy}) &  \\
English & a cow &  \\
\bottomrule
\end{tabularx}

\begin{grammarlens}[title={What you've built}]
The first of five animals.
\end{grammarlens}

\subsection*{Your turn}
\begin{itemize}
\raggedright
  \item \emph{Say it:} \emph{pasu}
  \item \emph{Say it:} it once more, slowly
  \item \emph{Say it:} \emph{pasu}, then \emph{kirāmam}, and put the one inside the other
  \item \emph{Recall:} say \emph{siṟitu}
  \item \emph{Read it:} \textbf{\ta{காற்று}}
\end{itemize}

\begin{tcolorbox}[breakable,colback=teal!4,colframe=teal!35!black,title={Before you move on}]
What does \ta{பசு} mean? (\textquotedblleft{}a cow\textquotedblright{}.) Say it once more.
\end{tcolorbox}
\section[\texorpdfstring{\ta{ஆடு} (āṭu)}{ஆடு (āṭu)}]{\ta{ஆடு} (āṭu) — a goat}

\label{lesson:TA-C58-goat}



\begin{quote}
Before the new one: what did \emph{pasu} mean?
\end{quote}

\subsection*{You'll want to know: \ta{ஆடு}}
\textbf{\ta{ஆடு}} (\emph{āṭu}) — \textquotedblleft{}a goat\textquotedblright{}.

Smaller, cheaper, and kept by houses with no land for a \ta{பசு}. It eats what it finds and needs nobody to cut for it.

Start it long — \ta{ஆ}, held — and put the \ta{ட} of \ta{கூடை} in the middle. That keeps it clear of \ta{அது}, which starts short and is a different word entirely. Kannada parts company here and says \ka{ಮೇಕೆ} (\emph{mēke}).

\textbf{In the family, and next door}

\noindent
\begin{tabularx}{\linewidth}{@{}>{\raggedright\arraybackslash}X>{\raggedright\arraybackslash}X>{\raggedright\arraybackslash}X@{}}
\toprule
\textbf{} & \textbf{word} & \textbf{same root} \\
\midrule
Kannada & \ka{ಮೇಕೆ} (\emph{mēke}) &  \\
Telugu & \te{మేక} (\emph{mēka}) &  \\
Malayalam & \ml{ആട്} (\emph{āṭŭ}) & yes \\
Hindi (neighbour) & \dv{बकरी} (\emph{bakrī}) &  \\
English & a goat &  \\
\bottomrule
\end{tabularx}

\begin{grammarlens}[title={What you've built}]
Two.
\end{grammarlens}

\subsection*{Your turn}
\begin{itemize}
\raggedright
  \item \emph{Say it:} \emph{āṭu}
  \item \emph{Say it:} it once more, slowly
  \item \emph{Say it:} \emph{pasu}, then \emph{āṭu}, and say which one a house keeps when it has no land
  \item \emph{Read it:} \textbf{\ta{நிறைய}}
  \item \emph{Recall:} say \emph{maṇal}
\end{itemize}

\begin{tcolorbox}[breakable,colback=teal!4,colframe=teal!35!black,title={Before you move on}]
What does \ta{ஆடு} mean? (\textquotedblleft{}a goat\textquotedblright{}.) Say it once more.
\end{tcolorbox}
\section[\texorpdfstring{\ta{கோழி} (kōḻi)}{கோழி (kōḻi)}]{\ta{கோழி} (kōḻi) — a hen}

\label{lesson:TA-C58-hen}



\begin{quote}
Before the new one: what did \emph{āṭu} mean?
\end{quote}

\subsection*{You'll want to know: \ta{கோழி}}
\textbf{\ta{கோழி}} (\emph{kōḻi}) — \textquotedblleft{}a hen\textquotedblright{}.

Hen or cock alike; the word does not split them, and where the difference matters a second word goes in front. Kannada \ka{ಕೋಳಿ} (\emph{kōḷi}) and Telugu \te{కోడి} (\emph{kōḍi}) are the same word three ways, and the middle letter is where they separate: Tamil's \ta{ழ}, Kannada's \ka{ಳ}, Telugu's \te{డ}.

It is the clock of a \ta{கிராமம்}. A \ta{கோழி} has a house awake well before anybody there needs to know the \ta{மணி}.

\textbf{In the family, and next door}

\noindent
\begin{tabularx}{\linewidth}{@{}>{\raggedright\arraybackslash}X>{\raggedright\arraybackslash}X>{\raggedright\arraybackslash}X@{}}
\toprule
\textbf{} & \textbf{word} & \textbf{same root} \\
\midrule
Kannada & \ka{ಕೋಳಿ} (\emph{kōḷi}) & yes \\
Telugu & \te{కోడి} (\emph{kōḍi}) & yes \\
Malayalam & \ml{കോഴി} (\emph{kōḻi}) & yes \\
Hindi (neighbour) & \dv{मुर्गी} (\emph{murgī}) &  \\
English & a hen &  \\
\bottomrule
\end{tabularx}

\begin{grammarlens}[title={What you've built}]
Three.
\end{grammarlens}

\subsection*{Your turn}
\begin{itemize}
\raggedright
  \item \emph{Say it:} \emph{kōḻi}
  \item \emph{Say it:} it once more, slowly
  \item \emph{Say it:} \emph{kōḻi}, then \emph{maṇi}, and say which one comes first in the morning
  \item \emph{Recall:} say \emph{kuṟaivu}
  \item \emph{Read it:} \textbf{\ta{சேறு}}
\end{itemize}

\begin{tcolorbox}[breakable,colback=teal!4,colframe=teal!35!black,title={Before you move on}]
What does \ta{கோழி} mean? (\textquotedblleft{}a hen\textquotedblright{}.) Say it once more.
\end{tcolorbox}
\section[\texorpdfstring{\ta{காக்கை} (kākkai)}{காக்கை (kākkai)}]{\ta{காக்கை} (kākkai) — a crow}

\label{lesson:TA-C58-crow}



\begin{quote}
Before the new one: what did \emph{kōḻi} mean?
\end{quote}

\subsection*{You'll want to know: \ta{காக்கை}}
\textbf{\ta{காக்கை}} (\emph{kākkai}) — \textquotedblleft{}a crow\textquotedblright{}.

Named off the noise it makes, and the sisters landed on the same trick from the same starting point: Kannada \ka{ಕಾಗೆ} (\emph{kāge}). Say it once and you have the etymology.

Of everything with wings around a Tamil house this is the one the house deals with daily. \ta{சாதம்} set out on the \ta{வாசல்} step is put there for the \ta{காக்கை}, and it comes down within the minute.

\textbf{In the family, and next door}

\noindent
\begin{tabularx}{\linewidth}{@{}>{\raggedright\arraybackslash}X>{\raggedright\arraybackslash}X>{\raggedright\arraybackslash}X@{}}
\toprule
\textbf{} & \textbf{word} & \textbf{same root} \\
\midrule
Kannada & \ka{ಕಾಗೆ} (\emph{kāge}) & yes \\
Telugu & \te{కాకి} (\emph{kāki}) & yes \\
Malayalam & \ml{കാക്ക} (\emph{kākka}) & yes \\
Hindi (neighbour) & \dv{कौआ} (\emph{kauā}) &  \\
English & a crow &  \\
\bottomrule
\end{tabularx}

\begin{grammarlens}[title={What you've built}]
Four.
\end{grammarlens}

\subsection*{Your turn}
\begin{itemize}
\raggedright
  \item \emph{Say it:} \emph{kākkai}
  \item \emph{Say it:} it once more, slowly
  \item \emph{Say it:} \emph{kākkai}, then \emph{kōḻi}, and say which one is kept and which one turns up
  \item \emph{Read it:} \textbf{\ta{தேவையில்லை}}
  \item \emph{Recall:} say \emph{pukai}
\end{itemize}

\begin{tcolorbox}[breakable,colback=teal!4,colframe=teal!35!black,title={Before you move on}]
What does \ta{காக்கை} mean? (\textquotedblleft{}a crow\textquotedblright{}.) Say it once more.
\end{tcolorbox}
\section[\texorpdfstring{\ta{பறவை} (paṟavai)}{பறவை (paṟavai)}]{\ta{பறவை} (paṟavai) — a bird}

\label{lesson:TA-C58-bird}



\begin{quote}
Before the new one: what did \emph{kākkai} mean?
\end{quote}

\subsection*{You'll want to know: \ta{பறவை}}
\textbf{\ta{பறவை}} (\emph{paṟavai}) — \textquotedblleft{}a bird\textquotedblright{}.

The wide word. Whatever flies is a \ta{பறவை}, the \ta{கோழி} and the \ta{காக்கை} taken in along with everything that was never given a name of its own.

Here the three go three ways: Kannada \ka{ಹಕ್ಕಿ} (\emph{hakki}), Telugu \te{పక్షి} (\emph{pakṣi}) off the Sanskrit, and Tamil's own \ta{பறவை}. The \ta{ற} in the middle is the hard one, the same \ta{ற} standing in \ta{ஆறு}.

\textbf{In the family, and next door}

\noindent
\begin{tabularx}{\linewidth}{@{}>{\raggedright\arraybackslash}X>{\raggedright\arraybackslash}X@{}}
\toprule
\textbf{} & \textbf{word} \\
\midrule
Kannada & \ka{ಹಕ್ಕಿ} (\emph{hakki}) \\
Telugu & \te{పక్షి} (\emph{pakṣi}) \\
Malayalam & \ml{പക്ഷി} (\emph{pakṣi}) \\
Hindi (neighbour) & \dv{चिड़िया} (\emph{chiṛiyā}) \\
English & a bird \\
\bottomrule
\end{tabularx}

\begin{grammarlens}[title={What you've built}]
Five, and the run is closed: a cow, a goat, a hen, a crow, and the word for everything that flies.
\end{grammarlens}

\subsection*{Your turn}
\begin{itemize}
\raggedright
  \item \emph{Say it:} \emph{paṟavai}
  \item \emph{Say it:} it once more, slowly
  \item \emph{Say it:} all five in order — \emph{pasu}, \emph{āṭu}, \emph{kōḻi}, \emph{kākkai}, \emph{paṟavai}
  \item \emph{Recall:} say \emph{ākaṭṭum}
  \item \emph{Read it:} \textbf{\ta{கனல்}}
\end{itemize}

\begin{tcolorbox}[breakable,colback=teal!4,colframe=teal!35!black,title={Before you move on}]
What does \ta{பறவை} mean? (\textquotedblleft{}a bird\textquotedblright{}.) Say it once more.
\end{tcolorbox}
